\section{Overview}
\label{sec:phase2}

Write a technical review paper on \textbf{Knowledge Graph Construction for Hindi / Indic Languages}, modeled after the Arabic KG survey (Ahmed et al., 2022). This will be a supervised undergraduate publication with guidance from the course professor.

\section{Why a Review Paper?}

\begin{itemize}[nosep]
    \item No novel implementation required --- pure reading, synthesis, and writing
    \item Follow the exact structure of the Arabic KG survey (Ahmed et al., 2022) as a template
    \item Manageable scope: Indic KG literature is smaller than Arabic, requiring $\sim$40 papers to read
    \item Builds academic writing skills for future publications
    \item Suitable for a first undergraduate paper regardless of journal quartile (Q1--Q4)
\end{itemize}

\section{Template Paper}

The paper \textit{``Arabic Knowledge Graph Construction: A close look in the present and into the future''} (Ahmed et al., 2022, Journal of King Saud University -- Computer and Information Sciences) provides the template structure:

\begin{enumerate}[nosep]
    \item \textbf{Introduction} --- Motivation for KGs, why Indic languages need dedicated KG research
    \item \textbf{Indic Languages Overview} --- Linguistic challenges: morphology, script diversity, code-mixing, resource scarcity
    \item \textbf{Knowledge Graph Overview} --- Definition, construction phases (NER, RE, entity linking, fusion), applications (QA, recommendation)
    \item \textbf{Existing Approaches} --- Survey of existing KG systems for Indic languages (Hindi, Urdu, Bengali, Tamil, etc.) and cross-lingual approaches
    \item \textbf{Challenges \& Problems} --- Data sparsity, tool availability, annotation quality, evaluation benchmarks
    \item \textbf{Future Directions} --- LLM-based KG construction, cross-lingual alignment, Indic Wikipedia mining
    \item \textbf{Conclusion}
\end{enumerate}

\section{Research Methodology}

\begin{enumerate}[nosep]
    \item Search major venues (ACL, EMNLP, NAACL, COLING, LREC, TALLIP) for Indic KG papers using keywords: ``Hindi knowledge graph,'' ``Indic knowledge graph construction,'' ``Urdu KG,'' ``Indo-Aryan KG''
    \item Categorize papers by: language covered, approach (rule-based / ML / hybrid), construction phase (NER / RE / entity linking / fusion), application domain
    \item Identify gaps and open problems (e.g., no large-scale Hindi KG benchmark, limited Urdu Wikidata coverage)
    \item Synthesize into a taxonomy table (similar to Ahmed et al.\ Fig.\ 4)
\end{enumerate}

\section{Expected Timeline}

\begin{table}[h]
\centering
\begin{tabularx}{\textwidth}{cX}
\toprule
\textbf{Week} & \textbf{Activity} \\
\midrule
4 & Read Arabic KG survey as template; search and collect 40+ Indic KG papers \\
5 & Read and annotate papers; categorize by approach and language \\
6 & Write first draft following template structure \\
7 & Iterate with professor feedback; revise \\
8 & Polish formatting, finalize references, submit \\
\bottomrule
\end{tabularx}
\end{table}

\section{Target Venues}

Suitable venues for an undergraduate survey paper:
\begin{itemize}[nosep]
    \item ACL Student Research Workshop (ACL-SRW)
    \item Workshop on NLP for Similar Languages, Varieties and Dialects (VarDial)
    \item Conference on Empirical Methods in NLP Student Workshop
    \item Journal of King Saud University -- Computer and Information Sciences (same venue as template)
    \item Any Q3/Q4 journal recommended by professor
\end{itemize}
